where
\[
 U(x)=\begin{pmatrix}1&x\\0&1\end{pmatrix},\quad
 D(u)=\begin{pmatrix}u&0\\0&u^{-1}\end{pmatrix},\quad
 S=\begin{pmatrix}0&-1\\1&0\end{pmatrix}.
\]
The rescaled matrix has determinant one.  For every odd modulus $c$, the element
$-2$ is a unit in $\Z/c\Z$, so the factorisation applies whenever $d$ is also a
unit modulo $c$.
\end{proposition}

\begin{proof}
Direct multiplication of the right-hand side gives the left-hand side after
using $\lambda(AB-kd)=1$ and $dd^{-1}=1$.  The determinant is
$\lambda(AB-kd)=1$.  If $c$ is odd, $(2,c)=1$, hence $-2$ is invertible modulo
$c$.
\end{proof}


\section{The abstract Hilbert--HMRD theorem and the physical-caller no-go}\label{sec:hmrd}

This section records the finite Hilbert-valued reciprocal estimate that survives
formal verification.  Its validity must be separated from the failure of the
particular physical caller for which it was originally proposed.

Let $s\ge1$, let $H$ be a finite-dimensional complex Hilbert space, and let
$D_0,W_0$ be finite index sets.  For coefficients $\alpha_d,\beta_w\in\C$,
vectors $u_d,v_w\in H$, residue classes $x_d,y_w\in\Z/s\Z$, and a phase
$f:\Z/s\Z\to\C$, define
\begin{equation}\label{eq:hmrd-packet}
 S_H=\sum_{d\in D_0}\sum_{w\in W_0}
       \alpha_d\beta_w f(x_dy_w)\ip{u_d}{v_w},
\end{equation}
where the pairing is linear in the first variable.  Put
\begin{align}
 B_1&=\sum_{w\in W_0}|\beta_w|\norm{v_w},\label{eq:B1}\\
 A_2^2&=\sum_{d\in D_0}|\alpha_d|^2\norm{u_d}^2,
 &B_2^2&=\sum_{w\in W_0}|\beta_w|^2\norm{v_w}^2.\label{eq:A2B2}
\end{align}

\classification{\newderivation}{Kernel checked in \path{Gate1B/HilbertHMRD/Fourier.lean}, \path{OperatorStability.lean}, and \path{Core.lean}.}
\begin{theorem}[Finite abstract Hilbert--HMRD]\label{thm:hmrd}
Assume that $f$ is supported on the units modulo $s$ and $|f(x)|\le1$.  Suppose
that
\begin{equation}\label{eq:hmrd-P}
 \norm{\sum_{d\in D_0}\alpha_d\eord{s}{m x_d}u_d}\le P
       \qquad (m\bmod s).
\end{equation}
Then the small-modulus branch is
\begin{equation}\label{eq:hmrd-small}
                 |S_H|\le\sqrt{\phi(s)}\,P B_1.
\end{equation}

Suppose, alternatively and compatibly, that $A\in(\Z/s\Z)^\times$ and there
are residue maps $p_d,q_w$ for which
\begin{equation}\label{eq:hmrd-phase}
              f(x_dy_w)=\eord{s}{A p_dq_w}.
\end{equation}
If every $p$-fibre has cardinality at most $c_D$ and every $q$-fibre has
cardinality at most $c_W$, then the energy branch is
\begin{equation}\label{eq:hmrd-energy}
             |S_H|\le\sqrt{s c_Dc_W}\,A_2B_2.
\end{equation}
Consequently,
\begin{equation}\label{eq:hmrd-min}
 |S_H|\le
 \min\left\{\sqrt{\phi(s)}\,PB_1,
             \sqrt{s c_Dc_W}\,A_2B_2\right\}.
\end{equation}
For sources lying in intervals of lengths $D'$ and $W'$, one may take
\[
              c_D=1+\lfloor D'/s\rfloor,
 \qquad       c_W=1+\lfloor W'/s\rfloor.
\]
\end{theorem}

\begin{proof}
Use the normalised finite Fourier coefficient
\[
 \widehat f(m)=\frac1s\sum_{x\bmod s}f(x)\eord{s}{-mx}.
\]
Fourier inversion gives
$f(x)=\sum_{m\bmod s}\widehat f(m)\eord{s}{mx}$, while Parseval and the unit
support give
\[
 \sum_{m\bmod s}|\widehat f(m)|^2
   =\frac1s\sum_{x\bmod s}|f(x)|^2
   \le\frac{\phi(s)}s.
\]
Cauchy--Schwarz therefore yields the exact Fourier $\ell^1$ cost
\begin{equation}\label{eq:fourier-l1}
              \sum_{m\bmod s}|\widehat f(m)|\le\sqrt{\phi(s)}.
\end{equation}
Expanding $f(x_dy_w)$ and grouping first in $d$ expresses the $m$-slice of
$S_H$ as a sum over $w$ of a pairing between $v_w$ and a twisted $d$-vector.
By \cref{eq:hmrd-P}, Cauchy--Schwarz in the Hilbert space, and
\cref{eq:B1}, the norm of every slice is at most $PB_1$.  Summing against
$|\widehat f(m)|$ and applying \cref{eq:fourier-l1} proves
\cref{eq:hmrd-small}.

For the energy branch, group the $d$-vectors and $w$-vectors by their residue
classes:
\[
 X_p=\sum_{\substack{d\in D_0\\p_d=p}}\alpha_du_d,
 \qquad
 Y_q=\sum_{\substack{w\in W_0\\q_w=q}}\beta_wv_w.
\]
Then \cref{eq:hmrd-phase} transforms the packet into
\[
       \sum_{p,q\bmod s}\eord{s}{Apq}\ip{X_p}{Y_q}.
\]
The scalar kernel $K(p,q)=\eord{s}{Apq}$ satisfies
\[
 \sum_{p\bmod s}K(p,q)\overline{K(p,q')}
  =\begin{cases}s,&q=q',\\0,&q\ne q',\end{cases}
\]
because $A$ is a unit.  Hence its scalar $\ell^2$ operator norm is $\sqrt s$.
Tensoring with the identity on a finite-dimensional Hilbert space preserves
that norm; the formal development proves both directions of this operator-norm
identity.  Therefore
\begin{equation}\label{eq:kernel-H-bound}
 \left|\sum_{p,q}\eord{s}{Apq}\ip{X_p}{Y_q}\right|
 \le\sqrt s
      \left(\sum_p\norm{X_p}^2\right)^{1/2}
      \left(\sum_q\norm{Y_q}^2\right)^{1/2}.
\end{equation}
Within a fibre of size at most $c_D$, Cauchy--Schwarz gives
\[
 \norm{X_p}^2\le c_D
   \sum_{d:p_d=p}|\alpha_d|^2\norm{u_d}^2.
\]
Summing over $p$ gives $\sum_p\norm{X_p}^2\le c_DA_2^2$; similarly
$\sum_q\norm{Y_q}^2\le c_WB_2^2$.  Substitution into
\cref{eq:kernel-H-bound} proves \cref{eq:hmrd-energy}, and taking the minimum
of two simultaneously valid upper bounds proves \cref{eq:hmrd-min}.  The
interval fibre count is at most one plus interval length divided by the modulus.
\end{proof}

\begin{remark}[Effective modulus]
The Hilbert theorem receives $s$ after conductor reduction.  In applications
where a phase initially has modulus $r$ and coefficient $A$, the effective
modulus $s=r/(A,r)$ must be produced before calling \cref{thm:hmrd}.  That
separate algebraic reduction is not hidden inside the theorem.
\end{remark}

\classification{\newreduction}{Logical no-go kernel checked in \path{Gate1B/R11/StatusBank.lean}; analytic nonclosure independently retained.}
\begin{proposition}[Physical-caller firewall]\label{prop:physical-caller}
The abstract theorem \cref{thm:hmrd} does not imply a bound for the physical
$R_{11}$ source unless all of the following are supplied:
\begin{enumerate}[label=\textup{(\alph*)},leftmargin=2.7em]
\item an exact factorisation of the physical delta-constrained source as a
      finite packet of the form \cref{eq:hmrd-packet};
\item the corresponding physical vector-valued M\"obius estimate required by
      the small branch or the full residue-energy dictionary required by the
      energy branch;
\item a global mode ledger proving that the packets so produced exhaust the
      physical source with the stated normalisation.
\end{enumerate}
No such data are constructed in the formal bank.  The previous
``physical delta plus single Poisson'' caller is therefore retired as
nonclosing.  This retirement does not invalidate \cref{thm:hmrd}.
\end{proposition}

\begin{proof}
In the formal layer, \texttt{PhysicalHilbertCaller},
\texttt{VectorMobiusPhysicalBound}, and \texttt{GlobalModeLedger} are explicit
propositions or structures that are never inhabited.  The only transport
result has the logical form ``if a physical caller is supplied, then the
abstract HMRD bound applies''.  Thus no unconditional physical conclusion can
be extracted.  Analytically, a single Poisson opening by itself neither proves
that the target tensor factorises as \cref{eq:hmrd-packet} nor supplies the
source norms in (b) and the exhaustive reassembly in (c).
\end{proof}

\section{The MHHC two-copy reduction}\label{sec:mhhc}

The first-moment row on one affine line is schematically
\begin{equation}\label{eq:first-moment-row}
        \sum_{t\in\cI_{A,k,C}}\mu(d_t)b^\perp_{A,k,C}(t).
\end{equation}
A Cauchy--Schwarz or dispersion step creates the source-exact two-copy
correlation
\begin{equation}\label{eq:mhhc-two-copy}
 \cO^\perp=
 \sum_{C}\sum_{A,k}\sum_{0<|h|<T_C}\sum_t
 \mu(d_t)\mu(d_t+Ah)
 b^\perp_{A,k,C}(t)\overline{b^\perp_{A,k,C}(t+h)}.
\end{equation}
The determinant identity \cref{eq:det2h} is preserved term by term.

\classification{\newreduction}{Exact determinant identity kernel checked; the reduction from the supplied physical row is inherited and not re-formalised.}
\begin{proposition}[Matched-Hecke two-copy geometry]\label{prop:mhhc-geometry}
For each term of \cref{eq:mhhc-two-copy},
\begin{equation}\label{eq:mhhc-det}
 d_tB_{t+h}-d_{t+h}B_t=2h,
\end{equation}
and
\begin{equation}\label{eq:matrix-hecke}
 \begin{pmatrix}d_t&d_{t+h}\\B_t&B_{t+h}\end{pmatrix}
 =
 \begin{pmatrix}d_t&A\\B_t&k\end{pmatrix}
 \begin{pmatrix}1&1\\0&h\end{pmatrix}.
\end{equation}
The first factor has determinant $2$, and the second has determinant $h$.
\end{proposition}

\begin{proof}
Equation \cref{eq:mhhc-det} is \cref{eq:det2h}.  Since
$(d_{t+h},B_{t+h})=(d_t,B_t)+h(A,k)$, multiplication gives
\cref{eq:matrix-hecke}.  The determinant statements follow from
$kd_t-AB_t=2$.
\end{proof}

For $h>0$, after determinant normalisation one obtains
\begin{equation}\label{eq:hecke-cell}
 (2h)^{-1/2}
 \begin{pmatrix}d_t&d_{t+h}\\B_t&B_{t+h}\end{pmatrix}
 =g_tu_{h,1},
\end{equation}
where $g_t\in\mathrm{SL}_2(\R)$ and
\[
 u_{h,1}=h^{-1/2}\begin{pmatrix}1&1\\0&h\end{pmatrix}.
\]
This is one cell of a Hecke correspondence, not the full Hecke operator.  The
exact completion of that cell in the upper-right residue is
\begin{equation}\label{eq:hecke-completion}
 f(u_{h,1})=\frac1h\sum_{\xi\bmod h}\eord{h}{-\xi}
                  \sum_{b\bmod h}\eord{h}{\xi b}f(u_{h,b}).
\end{equation}
It leaves both nonzero $b$-frequencies and the other factorisations $ad=h$.

\begin{warning}[One-copy and two-copy roles]
The rough-conductor theorem below is applied to the complete \emph{first
moment} before a fresh two-copy expansion.  It does not prove
\cref{eq:mhhc-two-copy} small on the remaining high frequencies.  After the low
projection is removed linearly, any later MHHC argument must be run on the
actual high-pass target $(I-P_C^{\mathrm{lo}})b^\perp$.
\end{warning}

\section{Mean-zero and Fourier structure}\label{sec:fourier-physical}

Choose an integer $Q_{A,k,C}$ satisfying
\begin{equation}\label{eq:Q-range}
       2T_C\le Q_{A,k,C}<4T_C
\end{equation}
(or any comparable cyclic length at least twice the physical interval), and
extend $b^\perp$ by zero to $\Z/Q\Z$.  We use the normalised coefficients
\begin{equation}\label{eq:cxi}
 c_{A,k,C}(\xi)=\frac1Q\sum_{t\bmod Q}b^\perp(t)\eord{Q}{-\xi t}.
\end{equation}
Then
\begin{equation}\label{eq:physical-fourier}
 c(0)=0,\qquad
 b^\perp(t)=\sum_{\xi\bmod Q}c(\xi)\eord{Q}{\xi t},
 \qquad
 \sum_{\xi\bmod Q}|c(\xi)|^2
  =\frac1Q\sum_t|b^\perp(t)|^2.
\end{equation}
In particular, if $|b^\perp(t)|\le B_0$, then
\begin{equation}\label{eq:c-energy}
